\documentclass[11pt]{article}
\usepackage[margin=0.82in]{geometry}
\usepackage{amsmath,amssymb,amsthm,mathtools}
\usepackage{graphicx}
\usepackage{microtype}
\usepackage[colorlinks=true,linkcolor=blue,urlcolor=blue,citecolor=blue]{hyperref}

\newtheorem{theorem}{Theorem}
\newtheorem{lemma}{Lemma}
\newcommand{\Ball}{\operatorname{Ball}}
\newcommand{\oa}{\operatorname{oa}}

\title{A peak majorant need only commute with the peak projection}
\author{Candidate full affirmative answer to a question in arXiv:1206.4022}
\date{13 August 2026}

\begin{document}
\maketitle

\noindent\fbox{\parbox{0.95\linewidth}{\textbf{Status: candidate full proof,
likely valid.}  Corollary 5.4 of Blecher--Read remains true when its
hypothesis that the positive invertible majorant $d$ commute with the peak
element $a$ is replaced by the weaker hypothesis $dq=qd$.}}

\section{Source question and exact scope}

The source is D.~P. Blecher and C.~J. Read, \emph{Operator algebras with
contractive approximate identities II}, arXiv:1206.4022v4, J. Funct. Anal.
264 (2013), 1049--1067 \cite{BRII}.  Corollary 5.4 assumes that $A$ is an
operator algebra in a unital $C^*$-algebra $B$, that $q$ is the peak of
$a\in\Ball(A^1)$, that $b\in A$ and $|b|$ commutes with $a$, and that
\[
 q b^*bq\le qdq
\]
for an invertible $d\in B_+$.  Under the additional condition $ad=da$, it
produces $x\in A$ with $xq=bq$ and $x^*x\le d$.

\begin{figure}[ht]
\centering
\includegraphics[width=0.94\linewidth]{figures/source_question_crop.png}
\caption{The end of Corollary 5.4 and Remark 2 on source PDF page 14.}
\end{figure}

The exact question is:
\begin{quote}
``We do not know if Corollary 5.4 is true if one replaces the condition that
$d$ commutes with $a$ with $d$ commuting with $q$.''
\end{quote}
We retain every other hypothesis of the corollary.  Products with $q$ below
are computed in $B^{**}$ via the canonical normal extension of the inclusion.

\section{Proof intuition}

Writing $f=d^{-1/2}$ turns $x^*x\le d$ into $\|xf\|\le1$.  Normalizing the
whole algebra is unhelpful because $f$ need not commute with $a$.  The key is
to normalize only the cyclic space generated by $b$ and $a$:
\[
 X=\overline{b\,\oa(1,a)\,f}.
\]
Let $J$ be the peak ideal in $\oa(1,a)$, consisting of the functions that
vanish at $q$, and put $Y=\overline{bJf}$.  The commutation assumptions imply
that every $q$-column from $X$ has range orthogonal to every $(1-q)$-column
from $X$.  Consequently right multiplication by $1-q$ is an $M$-projection
on $X^{**}$, with range $Y^{\perp\perp}$.  Thus $Y$ is an $M$-ideal and exact
best approximation is available.  The source inequality says precisely that
the quotient norm of $bf$ is at most one.

The fixed orthogonal range spaces in the proof are important: they ensure
that the max-norm decomposition survives passage to the bidual, avoiding an
invalid appeal to weak-star continuity of products.

\section{The cyclic weighted-space lemma}

We use the standard characterization that $Y$ is an $M$-ideal in a Banach
space $X$ when a projection $P:X^{**}\to Y^{\perp\perp}$ satisfies
\[
 \|z\|=\max\{\|Pz\|,\|z-Pz\|\}\qquad(z\in X^{**}).
\]
Every $M$-ideal is proximinal \cite{HWW}.

\begin{lemma}[Cyclic weighted peak ideal]\label{lem:cyclic}
Let $E$ be a unital commutative operator algebra contained in a unital
$C^*$-algebra $B$.  Let $q\in E^{**}\subset B^{**}$ be a central closed
projection and
\[
 J=\{h\in E:hq=0\},\qquad J^{\perp\perp}=E^{**}(1-q).
\]
Suppose $b\in B$, $f\in B_+$ is invertible, $fq=qf$, and $b^*b$ commutes
with $E$ and $q$.  Then, for
\[
 X=\overline{bEf},\qquad Y=\overline{bJf},
\]
$Y$ is an $M$-ideal in $X$.  Under the canonical realization of $X^{**}$ as
the weak-star closure of $X$ in $B^{**}$, its $M$-projection is
$Pz=z(1-q)$.  In particular,
\begin{equation}\label{eq:distance}
 \operatorname{dist}(z,Y)=\|zq\|\qquad(z\in X),
\end{equation}
and the distance is attained.
\end{lemma}

\begin{proof}
Represent $B^{**}$ normally and faithfully on a Hilbert space $H$.  If
$h,k\in E$, then
\begin{equation}\label{eq:mixed}
 (bhfq)^*\,bkf(1-q)=qfh^*b^*bkf(1-q)=0,
\end{equation}
because $q$ commutes with every factor between it and $1-q$.  By norm
density, the same holds with $bhf,bkf$ replaced by arbitrary elements of
$X$.

Let $H_q$ be the closed linear span of all ranges $xqH$, $x\in X$, and let
$H_0$ be the closed linear span of all ranges $x(1-q)H$, $x\in X$.
Equation \eqref{eq:mixed} gives $H_q\perp H_0$.  If a bounded net in $X$
converges weak-star to $z\in X^{\perp\perp}$, then its $q$-columns converge
weakly pointwise to $zq$ and its $(1-q)$-columns to $z(1-q)$.  Closed
subspaces of $H$ are weakly closed; hence
\[
 zqH\subset H_q,\qquad z(1-q)H\subset H_0.
\]
Thus $(zq)^*z(1-q)=0$, and therefore
\begin{equation}\label{eq:max}
 \|z\|=\max\{\|zq\|,\|z(1-q)\|\}
 \qquad(z\in X^{\perp\perp}).
\end{equation}

It remains to identify the range.  The hypothesis on $J^{\perp\perp}$ gives
a bounded net $(e_t)\subset J$ converging weak-star to $1-q$.  For $h\in E$,
using $fq=qf$,
\[
 bhf(1-q)=bh(1-q)f=\mathop{\mathrm{w^*\!\!-lim}}_t bhe_tf\in
 Y^{\perp\perp}.
\]
Norm density and then weak-star density show
$X^{\perp\perp}(1-q)\subset Y^{\perp\perp}$.  Conversely $Yq=0$, so right
multiplication by $1-q$ fixes $Y^{\perp\perp}$.  Hence
\[
 P(X^{\perp\perp})=X^{\perp\perp}(1-q)=Y^{\perp\perp}.
\]
Together with \eqref{eq:max}, this says that $P$ is an $M$-projection and
$Y$ is an $M$-ideal in $X$.  In the decomposition
$X^{**}=Y^{\perp\perp}\oplus_\infty X^{**}q$, the quotient norm of $z+Y$ is
$\|zq\|$, proving \eqref{eq:distance}.  Proximinality of $M$-ideals gives
attainment.
\end{proof}

\section{Affirmative answer}

\begin{theorem}\label{thm:answer}
Let $A\subset B$ be an operator algebra, where $B$ is a unital $C^*$-algebra,
and identify $A^1=A+\mathbb C1_B$.  Suppose $q\in(A^1)^{**}$ is the peak of
some $a\in\Ball(A^1)$, $b\in A$, and $|b|$ commutes with $a$.  If
$d\in B_+$ is invertible,
\[
 dq=qd,\qquad qb^*bq\le qdq,
\]
then there is $x\in A$ such that
\[
 xq=bq,\qquad x^*x\le d.
\]
Thus the answer to the question after Corollary 5.4 is yes.
\end{theorem}

\begin{proof}
Let $E=\oa(1,a)\subset A^1$.  Since $q$ is the peak of $a$, it belongs to
$E^{**}$, is closed there, and $1-q$ is the support projection of the ideal
\[
 J=\{h\in E:hq=0\};
\]
equivalently $J^{\perp\perp}=E^{**}(1-q)$ (this is the standard peak-ideal
correspondence; see \cite{Hay,BRII}).

Put $c=b^*b$ and $f=d^{-1/2}$.  The condition that $|b|$ commutes with $a$
implies that $c$ commutes with $E$.  Taking weak-star limits of $a^n$ also
shows that $c$ commutes with $q$.  The condition $dq=qd$ gives $fq=qf$.
Lemma \ref{lem:cyclic} therefore applies to
\[
 X=\overline{bEf},\qquad Y=\overline{bJf}.
\]

For $z_0=bf$, the compression hypothesis gives
\[
 \|z_0q\|^2
 =\|qfb^*bfq\|
 \le \|qfd f q\|=\|q\|=1.
\]
Here the inequality is conjugation of $qb^*bq\le qdq$ by $qfq$, and all
insertions of $q$ are legitimate because $q$ commutes with $c$ and $f$.
By \eqref{eq:distance} and proximinality, choose $y\in Y$ with
\[
 \|bf-y\|=\operatorname{dist}(bf,Y)=\|bfq\|\le1.
\]
Right multiplication by $f^{-1}=d^{1/2}$ maps $Y$ into
$\overline{bJ}\subset A$, because $bA^1\subset A$.  Define
\[
 x=b-yf^{-1}\in A.
\]
Since $Yq=0$ and $f^{-1}$ commutes with $q$, we have $xq=bq$.  Moreover
\[
 xf=bf-y,
\]
so $\|xf\|\le1$.  Consequently
$f x^*x f\le1$, and conjugating by $f^{-1}$ yields $x^*x\le f^{-2}=d$.
\end{proof}

\section{Verification, limitations, and novelty}

The proof was checked against four failure modes.  First, the source's
$2\times2$ example does not commute with $q$.  Second, the proof never claims
that $d$ or $f$ commutes with $a$.  Third, the bidual max decomposition is
proved using fixed orthogonal output subspaces, not joint weak-star
continuity of multiplication.  Fourth, both inclusions in
$X^{**}(1-q)=Y^{\perp\perp}$ are explicit.  The accompanying deterministic
check tests the mixed-product zero identity and max-norm decomposition in
finite matrix models; it is a sanity check, not a substitute for the proof.

The theorem answers exactly the question on source PDF page 14.  It retains
invertibility of $d$, the peak hypothesis, and commutation of $|b|$ with $a$;
it does not assert that any of those assumptions can be removed.

For novelty, bounded searches on 13 August 2026 covered the run indexes, the
exact question phrase, its main commutation variants, arXiv:1211.5010
\cite{Bnpi}, and arXiv:1407.1356 \cite{BRorder}.  No later primary source
answer was found, and the 2018 arXiv source version still contains the open
question.  Novelty confidence is moderate rather than definitive because the
argument uses standard $M$-ideal technology and may be implicit elsewhere.

\textbf{Human-review recommendation.}  Check Lemma \ref{lem:cyclic}, especially
the canonical realization of $X^{**}$ in $B^{**}$, persistence of the fixed
range orthogonality, and the range identity for the $M$-projection.  Then
check the final right-normalization calculation.  These are the only
substantive points.

\begin{thebibliography}{9}
\bibitem{BRII}
D.~P. Blecher and C.~J. Read,
\emph{Operator algebras with contractive approximate identities II},
J. Funct. Anal. \textbf{264} (2013), 1049--1067;
\url{https://arxiv.org/abs/1206.4022}.

\bibitem{Bnpi}
D.~P. Blecher,
\emph{Noncommutative peak interpolation revisited},
Bull. Lond. Math. Soc. \textbf{45} (2013), 1100--1106;
\url{https://arxiv.org/abs/1211.5010}.

\bibitem{BRorder}
D.~P. Blecher and C.~J. Read,
\emph{Order theory and interpolation in operator algebras},
Studia Math. \textbf{225} (2014), 177--228;
\url{https://arxiv.org/abs/1407.1356}.

\bibitem{Hay}
D.~M. Hay,
\emph{Closed projections and peak interpolation for operator algebras},
Integral Equations Operator Theory \textbf{57} (2007), 491--512;
\url{https://arxiv.org/abs/math/0512353}.

\bibitem{HWW}
P.~Harmand, D.~Werner, and W.~Werner,
\emph{$M$-ideals in Banach Spaces and Banach Algebras},
Lecture Notes in Mathematics 1547, Springer, 1993.
\end{thebibliography}

\end{document}
